\chapter{Operations on currents}\label{sec:current-operations}
Pushforwards, restrictions, and homotopies turn geometric constructions into identities with quantitative mass bounds. Throughout this section, $m$ is the current dimension and $N$ is the ambient Euclidean dimension. We retain the comass convention of Definition \ref{def:mass-current}. The main source is \cite{Sim18}, Chapter 6, Section 2, pp.~159--169; its initial Euclidean norm on forms must be distinguished from comass for nonrectifiable currents.

\begin{definition}[Support and normal currents]\label{def:normal-current}
The support $\operatorname{spt}T$ of a current in $U$ is the complement, relative to $U$, of the largest open set on which it vanishes on all test forms. A current is locally normal if $T$ and $\partial T$ have finite mass on every relatively compact open subset of $U$. It is normal with finite mass if $\mathbf M(T)+\mathbf M(\partial T)<\infty$. Normality is an analytic condition; it does not require integer multiplicities or rectifiability. Every locally integral current is locally normal. See \cite{Sim18}, Chapter 6, Sections 2--3, and \cite{Fed69}, Section 4.1.
\end{definition}

\begin{theorem}[Representation of currents with locally finite mass]\label{thm:current-mass-representation}
If $T$ has locally finite mass in $U\subset\mathbb R^N$, there are a Radon measure $\|T\|$ and a measurable $m$-vector field $\xi$ such that
\[
 T(\omega)=\int_U\langle\omega,\xi\rangle\,d\|T\|,
 \qquad |\xi|_{\mathrm{mass}}=1\quad\|T\|\text{-almost everywhere},
 \qquad \mathbf M_W(T)=\|T\|(W).
\]
Here $|\cdot|_{\mathrm{mass}}$ is the norm on $m$-vectors dual to comass. The field $\xi$ need not be simple. The current extends uniquely to compactly supported continuous forms, and its integral representation also defines evaluation on bounded Borel forms with compact support. This is the vector-measure representation theorem, obtained from Riesz representation and polar decomposition. Compare \cite{Sim18}, Chapter 6, equations (2.8)--(2.10), with the norm convention changed by duality as explained in Remark 2.7. A local proof is supplied below from the scalar representation theorem and finite-dimensional norm duality.
\end{theorem}
\begin{proof}
Lemma \ref{lem:current-continuous-extension} supplies the unique continuous extension. Apply Theorem \ref{thm:normed-vector-riesz} with the finite-dimensional space of $m$-forms equipped with comass. Its dual space is the space of $m$-vectors equipped with the mass norm, so its representation gives the required field and unit mass norm. Proposition \ref{prop:current-open-mass-variation} identifies the resulting measure with $\|T\|$. Integrating a bounded Borel form of compact support against the measurable polar field is well defined because that support has finite measure. This last operation defines the stated Borel evaluation; it is not a density assertion for the uniform norm on Borel forms.
\end{proof}

\begin{proposition}[Weak compactness from local mass bounds]\label{prop:current-weak-compactness}
Suppose $T_j$ are $m$-currents in $U$ and $\sup_j\mathbf M_W(T_j)<\infty$ for every $W\Subset U$. A subsequence converges weakly to an $m$-current with locally finite mass. If also $\sup_j\mathbf M_W(\partial T_j)<\infty$ for every such $W$, the limit is locally normal. This conclusion does not by itself preserve integer rectifiability; that is the additional content of Theorem \ref{thm:integral-compactness}. Compare \cite{Sim18}, Lemma 2.17, p.~162.
\end{proposition}
\begin{proof}
Extend all currents to continuous forms by Lemma \ref{lem:current-continuous-extension}. Take the countable scalar tests of Lemma \ref{lem:countable-radon-tests} and multiply them by the finitely many elements of a basis of $m$-forms. Their current values are bounded sequences, by the local mass bounds. Repeated subsequences and a diagonal extraction make these values converge simultaneously.

Every continuous compactly supported form is uniformly approximable by finite linear combinations of these tests with supports in a common compact set. On a relatively compact open neighborhood of that set the masses have a uniform bound. Thus the same uniform approximation and Cauchy estimate as in Lemma \ref{lem:measure-diagonal-functional} show convergence on every continuous compactly supported form. Write $L$ for the limit. Limits preserve linearity. For tests supported in any fixed relatively compact $W$, taking the limit of the original mass bound gives $|L(\omega)|\leq C_W\|\omega\|_{*,\infty}$. Restricting $L$ to smooth forms therefore defines a current of locally finite mass: convergence in the smooth test-form topology includes uniform convergence on a common compact support.

For each smooth test form, $\partial T_j(\eta)=T_j(d\eta)\to T(d\eta)=\partial T(\eta)$. If the boundary masses are locally uniformly bounded, taking suprema over the boundary test forms gives locally finite boundary mass for $T$. This proves normality without invoking a separate weak-star compactness theorem or asserting preservation of integer rectifiability.
\end{proof}

\begin{definition}[Restriction and contraction by a differential]\label{def:current-restriction}
For a locally finite-mass current and a Borel set $E\subset U$, define
\[
 (T\llcorner E)(\omega)=\int_E\langle\omega,\xi\rangle\,d\|T\|.
\]
More generally $T\llcorner u$ inserts a locally bounded Borel scalar function $u$ into this integral. For smooth $u$ and $m\geq1$, the $(m-1)$-current $T\llcorner du$ is defined by
\[
 (T\llcorner du)(\eta)=T(du\wedge\eta).
\]
The order $du\wedge\eta$ fixes the sign used below. Restricting a normal current to an arbitrary Borel set need not give a normal current. See \cite{Sim18}, Chapter 6, equations (2.11)--(2.12) and Section 4.
\end{definition}

\begin{proposition}[Boundary of a smooth scalar restriction]\label{prop:current-cutoff-boundary}
For a locally normal $m$-current, $m\geq1$, and smooth $u$,
\[
 \partial(T\llcorner u)=(\partial T)\llcorner u-T\llcorner du.
\]
In particular multiplication by a smooth compactly supported function produces a normal current of finite mass, with
\[
 \mathbf M\bigl(\partial(T\llcorner u)\bigr)
 \leq\int|u|\,d\|\partial T\|+\int|Du|\,d\|T\|.
\]
\end{proposition}
\begin{proof}
For a test $(m-1)$-form $\eta$, expand $d(u\eta)=du\wedge\eta+u\,d\eta$ and apply $T$. This gives the identity. If $\|\eta\|_*\leq1$, then $\|du\wedge\eta\|_*\leq|Du|$: on a unit simple $m$-vector this is the norm bound for contraction by $Du$. The mass estimates follow by integration and taking the supremum. Compare \cite{Sim18}, the computation on p.~170.
\end{proof}

\begin{definition}[Smooth pushforward and properness]\label{def:current-pushforward}
Let $f:U\to V$ be smooth and proper on $\operatorname{spt}T$: for every compact $K\subset V$, the set $f^{-1}(K)\cap\operatorname{spt}T$ is compact in $U$. Define
\[
 f_\#T(\omega)=T(\chi f^*\omega),
\]
where $\chi\in C_c^\infty(U)$ equals one on a neighborhood of $\operatorname{spt}T\cap f^{-1}(\operatorname{spt}\omega)$. Different choices give the same value because their difference vanishes near the part of the current being tested. In particular compact support of $T$ guarantees properness. See \cite{Sim18}, equation (2.29), p.~164.
\end{definition}

\begin{proposition}[Pushforward, boundary, and composition]\label{prop:current-pushforward-identities}
For smooth maps with the required properness,
\[
 \partial(f_\#T)=f_\#(\partial T),\qquad
 (g\circ f)_\#T=g_\#(f_\#T),\qquad
 \operatorname{spt}(f_\#T)\subset f(\operatorname{spt}T).
\]
The image on the right is relatively closed under the properness hypothesis. It can be strictly larger than the support of the pushforward because of cancellation.
\end{proposition}
\begin{proof}
The first two identities follow from $d(f^*\eta)=f^*(d\eta)$ and $(g\circ f)^*=f^*g^*$. In the boundary calculation, derivatives of the auxiliary cutoff vanish near $\operatorname{spt}T\cap f^{-1}(\operatorname{spt}\eta)$, and the remaining terms vanish where the test form does. To prove the support inclusion, a test form supported off $f(\operatorname{spt}T)$ pulls back to zero near the support of $T$. Properness makes that image closed: a convergent sequence of image points lies in a compact target neighborhood, whose inverse image on the support is compact. See \cite{Sim18}, equation (2.31).
\end{proof}

\begin{proposition}[Mass bound for a pushforward]\label{prop:current-pushforward-mass}
Let $T$ have locally finite mass and let $f$ be as in Definition \ref{def:current-pushforward}. For open $W\subset V$,
\[
 \mathbf M_W(f_\#T)\leq\int_{f^{-1}(W)}\|Df(x)\|^m\,d\|T\|(x).
\]
Thus $\mathbf M(f_\#T)\leq L^m\mathbf M(T)$ if $\|Df\|\leq L$ on $\operatorname{spt}T$. For $m=0$, use $L^0=1$, including when $L=0$. The constant is one with the comass convention; compare \cite{Sim18}, Remark 2.30(1), which also discusses the other norm convention.
\end{proposition}
\begin{proof}
On an orthonormal $m$-frame, the pullback evaluates $\omega$ on its image under $Df$, whose simple-vector norm is at most $\|Df\|^m$. Hence $\|f^*\omega\|_*\leq\|Df\|^m\|\omega\|_*$. Integrate this bound against the polar representation of $T$ and take the supremum over test forms supported in $W$. The image-vector estimate is dual to the mass norm, so it does not assume that the polar field of a general current is simple.
\end{proof}

\begin{theorem}[Lipschitz pushforwards of integral currents]\label{thm:lipschitz-current-pushforward}
Let $U\subset\mathbb R^N$ and $W\subset\mathbb R^q$ be open. For a locally integral $m$-current $T$ in $U$, $m\geq0$, and a locally Lipschitz map $f:U\to W$ proper on $\operatorname{spt}T$, the pushforward is well defined and locally integral. It commutes with boundary. If $f$ is globally $L$-Lipschitz, then
\[
 \mathbf M(f_\#T)\leq L^m\mathbf M(T),\qquad
 \mathbf M(\partial f_\#T)\leq L^{m-1}\mathbf M(\partial T)\quad(m\geq1).
\]
For $T=\llbracket M,\theta,\tau\rrbracket$, its action is the integral of the test form against $\Lambda^m D^Mf(\tau)\theta$ on $M$, where $D^Mf$ is the approximate tangential derivative. The signed multiplicities of overlapping images are added before taking mass. Sources: \cite{Sim18}, Lemma 2.36, pp.~166--167, and Lemma 3.12 with Remark 3.13, pp.~174--175. The analytic construction, its tangential identification, and preservation of integer multiplicity are now proved below. They require chart pullback convergence and signed image multiplicities; ambient almost-everywhere differentiability alone is insufficient.
\end{theorem}
\begin{proof}
Discard the Hausdorff-null part of the rectifiable carrier lying outside $\operatorname{spt}T$, which is permitted by its positive integer mass density. First construct the locally normal analytic pushforward by Lemma \ref{lem:local-proper-normal-pushforward}. A locally integral current is locally normal because both it and its boundary have locally finite mass. Fix a compact target set $K$ and choose the cutoff $\chi$ and global extension $F$ as in that lemma. Proposition \ref{prop:normal-geometric-pushforward-agreement} applies to $T\llcorner\chi$, whose multiplicity is real and whose underlying rectifiable set still has locally finite Hausdorff measure. Since $\chi=1$ near the part of the support mapping into $K$, its geometric formula on tests supported in $K$ is exactly
\[
 \int_M\theta(x)\langle\omega(f(x)),
       \Lambda^m d^Mf_x(\tau(x))\rangle\,d\mathcal H^m(x).
\]
For $m=0$ this is the ordinary pushforward of the locally finite signed atomic measure. Properness makes the number of source atoms over each compact target set finite.

The zero-dimensional case is therefore already locally integer rectifiable. For $m\geq1$, choose a compact source neighborhood $B\Subset U$ containing the support of $\chi$ and on a neighborhood of which the extension $F$ agrees with $f$. The restricted current $R=T\llcorner B$ is integer rectifiable and has finite mass; its boundary need not have finite mass, and no normality is used for $R$. Theorem \ref{thm:integer-rectifiable-geometric-image} makes $G_FR$ integer rectifiable. Its tangential formula agrees with that of the analytic pushforward on $K$: every source point contributing to a test supported in $K$ lies where $\chi=1$ and hence belongs to $B$. The same argument on a slightly larger compact target neighborhood gives a local integer rectifiable representation near $K$. The signed multiplicity formula in that theorem is the claimed addition of overlapping sheets. A countable target cover gives local integer rectifiability throughout $W$.

Boundary commutation holds by Lemma \ref{lem:local-proper-normal-pushforward}. Apply the preceding integer-image argument also to $\partial T$, whose support lies in $\operatorname{spt}T$ and inherits properness. Thus the boundary image is locally integer rectifiable as well, proving local integrality. If the target dimension is smaller than the current dimension, the corresponding exterior vector and current are zero; boundary commutation remains supplied by the analytic construction.

If $f$ is globally $L$-Lipschitz, its intrinsic differential has norm at most $L$. The geometric formula therefore gives $\|f_\#T\|(K)\leq L^m\|T\|(f^{-1}(K))$ for compact target sets, or the same estimate first on open neighborhoods and then by regularity. Exhaustion yields the stated total mass estimate, allowing infinite values. Apply the same argument to $\partial T$ for its bound. Use the convention $L^0=1$ in dimension zero. The proof uses the local area/tangent chain and the normal-current construction, without appealing to Federer--Fleming compactness or claiming that a smooth cutoff has integer multiplicity.
\end{proof}

\begin{example}[Projection can cancel two sheets]\label{ex:current-projection-cancellation}
Let $I_+$ and $I_-$ be the horizontal segments $[0,1]\times\{1\}$ and $[0,1]\times\{-1\}$, both oriented toward increasing first coordinate. Set $T=\llbracket I_+\rrbracket-\llbracket I_-\rrbracket$, and let $p(x,y)=x$. Then $\mathbf M(T)=2$ while $p_\#T=0$: both segment integrals become the same integral over $[0,1]$ with opposite signs. The associated unoriented varifold has pushforward of mass $2$ instead. This illustrates the signed-multiplicity distinction in \cite{Sim18}, Remark 3.13.
\end{example}

\begin{definition}[The interval product]\label{def:current-interval-product}
For a compactly supported finite-mass $m$-current $T$ in $\mathbb R^N$, write a smooth $(m+1)$-form on $\mathbb R\times\mathbb R^N$ as $dt\wedge\alpha(t,x)+\beta(t,x)$, with $\alpha$ and $\beta$ containing only spatial differentials. Define
\[
 (\llbracket[0,1]\rrbracket\times T)(dt\wedge\alpha+\beta)
 =\int_0^1 T(\alpha(t,\cdot))\,dt.
\]
The interval direction comes first in the orientation. The product has mass at most $\mathbf M(T)$ and support in $[0,1]\times\operatorname{spt}T$. For an integer rectifiable $T$, the product is integer rectifiable by the product representation of its oriented tangent planes. See \cite{Sim18}, equations (2.19)--(2.25) and Remark 3.2(2).
\end{definition}

\begin{proposition}[Boundary of the interval product]\label{prop:current-product-boundary}
For a compactly supported normal current $T$, let $i_t(x)=(t,x)$. Then
\[
 \partial(\llbracket[0,1]\rrbracket\times T)
 =(i_1)_\#T-(i_0)_\#T-\llbracket[0,1]\rrbracket\times\partial T.
\]
In particular the interval product of an integral current is integral. See \cite{Sim18}, equation (2.26), p.~164.
\end{proposition}
\begin{proof}
Write a test $m$-form as $dt\wedge\gamma+\delta$. The component of its exterior derivative containing $dt$ is $dt\wedge(\partial_t\delta-d_x\gamma)$. Integrating gives
$T(\delta(1,\cdot))-T(\delta(0,\cdot))-\int_0^1\partial T(\gamma(t,\cdot))\,dt$,
which is the stated formula. For dimension zero the term containing $\gamma$ is absent. The mass bounds for the product and its boundary give normality; the integer rectifiable product description gives integrality when $T$ is integral.
\end{proof}

\begin{proposition}[The homotopy formula]\label{prop:current-homotopy}
Let $T$ be a compactly supported normal current, and let $h:[0,1]\times\mathbb R^N\to\mathbb R^Q$ be smooth near $[0,1]\times\operatorname{spt}T$. Put $f=h(0,\cdot)$ and $g=h(1,\cdot)$. Then
\[
 g_\#T-f_\#T
 =\partial h_\#(\llbracket[0,1]\rrbracket\times T)
 +h_\#(\llbracket[0,1]\rrbracket\times\partial T).
\]
For maps defined only near this compact set, cutoffs localize the construction. See \cite{Sim18}, equation (2.32), p.~166.
\end{proposition}
\begin{proof}
Apply $h_\#$ to Proposition \ref{prop:current-product-boundary}, commute pushforward with boundary, and use $h\circ i_0=f$, $h\circ i_1=g$. Moving the product-boundary term to the other side yields the formula. Both terms matter when $T$ has nonzero boundary.
\end{proof}

\begin{definition}[A flat norm for real currents]\label{def:real-current-flat-norm}
For a compactly supported $m$-current in $\mathbb R^N$, define
\[
 \mathcal F_{\mathbb R}(T)=\inf_{T=R+\partial S}
 \bigl(\mathbf M(R)+\mathbf M(S)\bigr),
\]
where $R$ and $S$ are compactly supported real currents of finite mass of dimensions $m$ and $m+1$. The value is allowed to be infinite. This choice of competitors is different from the integer polyhedral competitors in Definition \ref{def:flat-norm}; for a polyhedral integer chain it gives a value no larger than that earlier infimum. No equality of these two infima is asserted here. Compare the flat topology in \cite{Fed69}, Section 4.1, and the homotopy estimates in \cite{Sim18}, Section 2.
\end{definition}

\begin{proposition}[Flat control of test forms and boundaries]\label{prop:real-flat-control}
For every compactly supported smooth $m$-form $\omega$,
\[
 |T(\omega)|\leq\mathcal F_{\mathbb R}(T)
 \max\{\|\omega\|_{*,\infty},\|d\omega\|_{*,\infty}\}.
\]
Consequently $\mathcal F_{\mathbb R}(T)=0$ implies $T=0$, flat convergence implies weak convergence, and
$\mathcal F_{\mathbb R}(\partial T)\leq\mathcal F_{\mathbb R}(T)$.
On the vector space where it is finite, $\mathcal F_{\mathbb R}$ is a norm.
\end{proposition}
\begin{proof}
For $T=R+\partial S$, write $T(\omega)=R(\omega)+S(d\omega)$ and use the two mass bounds. Take the infimum. If the infimum is zero, all test evaluations vanish. For the boundary estimate use $\partial T=\partial R$ as a decomposition with zero residual and filling $R$. Triangle inequality follows by adding decompositions; scaling decompositions gives absolute homogeneity. These arguments also explain why the test-form topology is weaker.
\end{proof}

\begin{proposition}[Nondegeneracy of the integer polyhedral flat metric]\label{prop:polyhedral-flat-nondegeneracy}
The flat norm in Definition \ref{def:flat-norm} is symmetric, subadditive, and zero only for the zero chain. Thus its metric completion does not identify distinct polyhedral chains.
\end{proposition}
\begin{proof}
Reversing the signs of a residual and a filling proves symmetry; adding two decompositions and using mass subadditivity proves the triangle inequality. For nondegeneracy, view each oriented polyhedron as its integration current. Stokes' formula identifies its polyhedral boundary with its current boundary. Thus every integer polyhedral decomposition is an admissible real-current decomposition, and
$\mathcal F_{\mathbb R}(P)\leq\mathcal F(P)$.
If $\mathcal F(P)=0$, Proposition \ref{prop:real-flat-control} shows that its current is zero. A nonzero subdivided polyhedral chain has a face with nonzero coefficient. Choose a point in its relative interior away from the other faces and a test form supported in a sufficiently small ball, with positive evaluation on its orientation. Its integral against the chain is nonzero, a contradiction. For dimension zero use a test function supported near a point with nonzero coefficient. This proves the metric assertion without identifying the real-current and integer-polyhedral flat infima.
\end{proof}

\begin{proposition}[A quantitative affine homotopy estimate]\label{prop:current-homotopy-estimate}
Let $m\geq1$, let $T$ be compactly supported and normal, and let smooth maps $f,g:\mathbb R^N\to\mathbb R^Q$ satisfy $\|Df\|,\|Dg\|\leq L$ near $\operatorname{spt}T$ and $|f-g|\leq\delta$ on that support. Then
\[
 \mathcal F_{\mathbb R}(g_\#T-f_\#T)
 \leq\delta\bigl(L^m\mathbf M(T)+L^{m-1}\mathbf M(\partial T)\bigr).
\]
For $m=1$, the second power is interpreted as one. This is the global finite-mass version of \cite{Sim18}, equation (2.34).
\end{proposition}
\begin{proof}
Use $h(t,x)=(1-t)f(x)+tg(x)$. The image of the interval direction is $g-f$, of norm at most $\delta$; each spatial differential has norm at most $L$. The product formula and the comass estimate therefore give
\[
 \mathbf M\bigl(h_\#(\llbracket[0,1]\rrbracket\times T)\bigr)
 \leq\delta L^m\mathbf M(T),
\]
and the same estimate with $\partial T$ gives $\delta L^{m-1}\mathbf M(\partial T)$. Insert these two currents as the filling and residual in the homotopy formula. For general polar vectors the estimate follows by duality with comass, just as in Proposition \ref{prop:current-pushforward-mass}.
\end{proof}

\begin{definition}[The cone over a compactly supported current]\label{def:current-cone}
For a compactly supported normal current $T$ in $\mathbb R^N$ and $a\in\mathbb R^N$, set
\[
 C_aT=h_\#(\llbracket[0,1]\rrbracket\times T),
 \qquad h(t,x)=a+t(x-a).
\]
This definition also applies to zero-dimensional currents. Its support lies in the union of the segments from $a$ to $\operatorname{spt}T$. See \cite{Sim18}, equations (2.37)--(2.39), p.~167.
\end{definition}

\begin{proposition}[Cone identity and mass estimate]\label{prop:current-cone-bound}
For $m\geq1$ and $T$ as above,
\[
 \partial C_aT=T-C_a(\partial T),\qquad
 \mathbf M(C_aT)\leq\frac{R}{m+1}\mathbf M(T)
 \quad\text{if }\operatorname{spt}T\subset\overline B_R(a).
\]
If $T$ is integral, then so is $C_aT$. In particular an integral cycle of positive dimension supported in a ball has an integral filling in the same ball, with the displayed mass bound.
\end{proposition}
\begin{proof}
In the homotopy formula, $g$ is the identity and $f$ is constant. A constant map sends every positive-dimensional current to zero, so the boundary identity follows. The differential sends the interval direction to $x-a$ and spatial $m$-vectors to $t^m$ times themselves. Thus the mass is at most
$\int_0^1 Rt^m\mathbf M(T)\,dt=R\mathbf M(T)/(m+1)$.
The interval product is integral when $T$ is, and its smooth pushforward remains integral by Theorem \ref{thm:lipschitz-current-pushforward}. Convexity of the ball controls the support.
\end{proof}

\begin{example}[The zero-dimensional augmentation condition]\label{ex:zero-current-augmentation}
For $T=\sum_{i=1}^k a_i\delta_{p_i}$ with integer coefficients, the cone is a sum of oriented segments and
\[
 \partial C_aT=T-\left(\sum_i a_i\right)\delta_a.
\]
Thus a zero-dimensional current has a compactly supported one-current filling when $\sum_i a_i=0$. This condition is also necessary: if $T=\partial S$, test against a compactly supported function identically one near $\operatorname{spt}S$ to get $T(1)=S(d1)=0$, with the cutoff understood. This explains why the positive-dimensional cone identity requires a separate statement in dimension zero.
\end{example}
